\refstepcounter{section}
\section*{Lecture 10: Connections are important}
\addcontentsline{toc}{section}{Lecture 10: Connections are important}
In the previous section, we had seen that the concept of being \textit{local} was introduced in order to approximate a frame where the effect of gravity is eliminated. Using such a motivation, we have:
\begin{principle}[Einstein's Equivalence Principle]
All freely falling frames under the influence of \textit{gravity} are \textit{locally inertial}, that is, Lorentz frames where special relativity is valid and there is absence of gravity
\end{principle}
Consider the case of the falling lift as before. Then, if we write the equation of the ball from the $O$ frame, we have:
$$\dv{\vb{r}}{t} = -\vb{g} \implies \vb{r} = \vb{r}_i - \frac{1}{2}\vb{g}t^2$$
where $\vb{r}_i$ is the initial position of the ball. Now, from the $O'$ frame, as no force acts on the ball (note that the observer sees the ball floating in front of them) we have:
$$\dv[2]{{\vb{r}'}}{{t'}} = 0\implies \vb{r}' = \vb{r}_i'$$
From the above equations we have:
$$\vb{r}' = \vb{r}_i'= (\vb{r}_i -\vb{r}_0) = \qty(\vb{r}+ \frac{1}{2}\vb{g}t^2) - \vb{r}_0 $$
This essentially describes a \textit{coordinate transformation} where $\vb{r}\rightarrow \vb{r}' $ and $t\rightarrow t'$. However, note that a term $t^2$ appears in the expression and hence, the coordinate transformations are \textit{non-linear}. In general, we can have the coordinate transformation matrices themselves to be dependent on the coordinates, that is, $\trmat{\mu}{\nu} \equiv \trmat{\mu}{\nu}(x^\alpha)$\\[0.2cm]
If we allow non-linear coordinate transformations, we can go to a frame where \textit{gravity is absent!}\\[0.2cm]
Thus, we essentially have a recipe to tackle gravity. We first go to a frame where gravity is absent (the equivalence principle guarantees the existence of such a frame), do our calculations in that frame, then transform back to our original coordinate system. \\[0.2cm]
Consider a particle moving freely under a purely gravitational force. Now, suppose $O'$ is such a frame where gravity is absent. Then we have:
$$\dv[2]{x'^\mu}{\tau}=0$$
How do we obtain the equation of motion in frame $O$? Well, we will just transform to that frame. Since $x'^\mu$ is dependent on the coordinates of the other frame $x^\nu$, we can write,
\begin{align*}
    0 &=\dv{\tau}\qty(\dv{x'^\mu}{\tau})\\
    &=\dv{\tau}\qty({\pdv{x'^\mu}{x^\nu}\dv{x^\nu}{\tau}} )\\
    &=\pdv{x'^\mu}{x^\nu} \dv[2]{x^\nu}{\tau}+ \dv{x^\nu}{\tau}\dv{\tau}\qty(\pdv{x'^\mu}{x^\nu})\\
    &=\pdv{x'^\mu}{x^\nu} \dv[2]{x^\nu}{\tau} + \pdv{x'^\mu}{x^\nu}{x^\alpha}\dv{x^\alpha}{\tau}\dv{x^\nu}{\tau}
\end{align*}
The first term seems to contain the acceleration in the $O$ frame and it would be nice to isolate this term. For that, let us multiply both sides by $\pdv{x^\beta}{x'^\mu}$ and contract using the fact that $\pdv{x^\beta}{x'^\mu}\pdv{x'^\mu}{x^\nu}=\tsrud{\delta}{\beta}{\nu}$
\begin{align*}
   0&= \pdv{x^\beta}{x'^\mu}\pdv{x'^\mu}{x^\nu} \dv[2]{x^\nu}{\tau} +  \pdv{x^\beta}{x'^\mu}\pdv{x'^\mu}{x^\nu}{x^\alpha}\dv{x^\alpha}{\tau}\dv{x^\nu}{\tau}\\
&=\tsrud{\delta}{\beta}{\nu}\dv[2]{x^\nu}{\tau}+ \pdv{x^\beta}{x'^\mu}\pdv{x'^\mu}{x^\nu}{x^\alpha}\dv{x^\alpha}{\tau}\dv{x^\nu}{\tau}\\
&=\dv[2]{x^\beta}{\tau}+ \pdv{x^\beta}{x'^\mu}\pdv{x'^\mu}{x^\nu}{x^\alpha}\dv{x^\alpha}{\tau}\dv{x^\nu}{\tau}\\
&=\dv[2]{x^\beta}{\tau} + \chris{\beta}{\nu}{\alpha} u^\alpha u^\nu
\end{align*}
Here we have defined the following:
$$\chris{\beta}{\nu}{\alpha} := \pdv{x^\beta}{x'^\mu}\pdv{x'^\mu}{x^\nu}{x^\alpha} = \itrmat{\beta}{\mu}\ \pdv{x^\nu} \qty(\trmat{\mu}{\alpha})$$
The above entity is termed as \textit{connection coefficients}. The name connection coefficients comes from the fact that they connect neighboring points and tell us how to calculate the rate of change of a vector field from one point to another nearby, even though the coordinate system may itself be changing. The equation of motion thus obtained is:
\begin{equation}\label{eqn:geod}
    \dv[2]{x^\beta}{\tau} = -\chris{\beta}{\nu}{\alpha} u^\alpha u^\nu
\end{equation}
The above equation is called the \textit{geodesic equation} and resembles some kind of \textit{pseudo-force} or \textit{non-inertial force}. Note that the connection coefficients in its present form, depends on both $\fvec{x}$ and $\fvec{x}'$ which is kind of problematic since we might not always have knowledge of both the frames. \\[0.2cm]
Since the primed frame is locally inertial, the metric is \textit{Minkowski}, that is,$$g'_{{\mu}{\nu}} = \pdv{x^\alpha}{x'^\mu}\pdv{x^\beta}{x'^\nu}g_{{\alpha}{\beta}}$$
Inverting the relation, we have:
$$g_{{\mu}{\nu}} = \pdv{x'^\alpha}{x^\mu}\pdv{x'^\beta}{x^\nu}\eta_{{\alpha}{\beta}}$$
From the definition of the affine connection, we can see that, since partial derivatives commute, the affine connections are \textit{symmetric} in the lower indices. Thus, 
$$\chris{\alpha}{\mu}{\nu} =\chris{\alpha}{\nu}{\mu}$$
Now, let us contract the upper index of the affine connection definition which will be used later:
\begin{align*}
    \pdv{x'^\lambda}{x^\beta}\chris{\beta}{\nu}{\alpha} = \pdv{x'^\lambda}{x^\beta}\pdv{x^\beta}{x'^\mu}\pdv{x'^\mu}{x^\nu}{x^\alpha} = \tsrud{\delta}{\lambda}{\mu}\pdv{x'^\mu}{x^\nu}{x^\alpha}=\pdv{x'^\lambda}{x^\nu}{x^\alpha} = \partial_\nu\qty(\pdv{x'^\lambda}{x^\alpha})
\end{align*}
Then, let us differentiate both sides with $x^\rho$ and using the above found relation, we obtain:
\begin{align*}
    \partial_\rho g_{\mu\nu} &= \partial_\rho \qty(\pdv{x'^\alpha}{x^\mu})\times \pdv{x'^\beta}{x^\nu}\eta_{{\alpha}{\beta}} + \partial_\rho \qty(\pdv{x'^\beta}{x^\nu})\times \pdv{x'^\alpha}{x^\mu}\eta_{{\alpha}{\beta}}\\
    &=\chris{\sigma}{\rho}{\mu}\pdv{x'^\alpha}{x^\sigma} \pdv{x'^\beta}{x^\nu}\eta_{{\alpha}{\beta}} +\chris{\sigma}{\rho}{\nu}\pdv{x'^\beta}{x^\sigma}\pdv{x'^\alpha}{x^\mu}\eta_{{\alpha}{\beta}}\\
    &=\chris{\sigma}{\rho}{\mu}g_{\sigma\nu} + \chris{\sigma}{\rho}{\nu}g_{\mu\sigma}
\end{align*}
We thus obtain:
\begin{equation}\label{eqn:chris1}
    {\partial_\rho g_{\mu\nu} = \chris{\sigma}{\rho}{\mu}g_{\sigma\nu} + \chris{\sigma}{\rho}{\nu}g_{\mu\sigma}}
\end{equation}
Since these are true for any indices, we obtain two more equations by the exchange $\rho \leftrightarrow \mu$ and $\rho \leftrightarrow \nu$:
\begin{align}
    \partial_\mu g_{\rho\nu} &= \chris{\sigma}{\mu}{\rho}g_{\sigma\nu} + \chris{\sigma}{\mu}{\nu}g_{\rho\sigma}\label{eqn:chris2}\\
    \partial_\nu g_{\mu\rho} &= \chris{\sigma}{\nu}{\mu}g_{\sigma\rho} + \chris{\sigma}{\nu}{\rho}g_{\mu\sigma}\label{eqn:chris3}
\end{align}
Adding Eq. \ref{eqn:chris2} and Eq. \ref{eqn:chris3} and then subtracting Eq. \ref{eqn:chris1}, we obtain:
\begin{alignat*}{3}
    &\partial_\mu g_{\rho\nu}+\partial_\nu g_{\mu\rho}-\partial_\rho g_{\mu\nu} &&=\mathcolor{red}{\chris{\sigma}{\mu}{\rho}g_{\sigma\nu}} + \mathcolor{blue}{\chris{\sigma}{\mu}{\nu}g_{\rho\sigma}} +\mathcolor{blue}{\chris{\sigma}{\nu}{\mu}g_{\sigma\rho}} + \mathcolor{OliveGreen}{\chris{\sigma}{\nu}{\rho}g_{\mu\sigma}}-\mathcolor{red}{\chris{\sigma}{\rho}{\mu}g_{\sigma\nu}}- \mathcolor{OliveGreen}{\chris{\sigma}{\rho}{\nu}g_{\mu\sigma}}\\
    &\implies\qquad\qquad\quad \chris{\sigma}{\mu}{\nu}g_{\rho\sigma} &&= \frac{1}{2} \qty(\partial_\mu g_{\rho\nu}+\partial_\nu g_{\mu\rho}-\partial_\rho g_{\mu\nu})\\
    &\implies\qquad\quad g^{\rho\lambda} \chris{\sigma}{\mu}{\nu}g_{\rho\sigma} &&=\frac{1}{2}g^{\rho\lambda} \qty(\partial_\mu g_{\rho\nu}+\partial_\nu g_{\mu\rho}-\partial_\rho g_{\mu\nu})\\
    &\implies \qquad\qquad\quad\tsrud{\delta}{\lambda}{\sigma} \chris{\sigma}{\mu}{\nu} &&=\frac{1}{2}g^{\rho\lambda} \qty(\partial_\mu g_{\rho\nu}+\partial_\nu g_{\mu\rho}-\partial_\rho g_{\mu\nu})
\end{alignat*}
The red and the green terms cancel first and then we contract the metric on the left to obtain an extremely important identity for the affine connection: 
\begin{equation}
    \chris{\lambda}{\mu}{\nu} =\frac{1}{2}g^{\rho\lambda} \qty(\partial_\mu g_{\rho\nu}+\partial_\nu g_{\mu\rho}-\partial_\rho g_{\mu\nu})
\end{equation}
We obtained an expression for the affine connection in terms of the metric for the space. In this form, these are called the \textit{Christoffel symbols} and denoted by: 
$$\left\{ \begin{matrix} \lambda \\ \mu  \nu \end{matrix}\right\}$$
NOTE: We can remember the equation like this: In the left hand side, the upper index is $\lambda$ which comes only in the metric in the right side. The derivatives with $+$ sign are with respect to the lower indices $\mu$ and $\nu$ and the index of the metric in the derivative is a dummy index $\rho$ and the other lower index not appearing in the derivative. Finally, the derivative with $-$ sign is with respect to the dummy index and the metric contains both the lower indices of the left.\\[0.2cm]
Hence, we see that the equation of motions are known provided the Christoffel symbols are known, which are known if metric is known! 
\subsection{Action tells a lot!}
We will now consider one of the most fundamental quantitites in physics, the \textit{action}, albeit in the \textit{relativistic} sense.
$$\Ss = \int \dd t\ L = \int \dd t \ (T-V)$$ Note that, action must be invariant under $\mathfrak{3}-$rotation in Newtonian mechanics. Since relativisitc case should also converge to the classical action in the Newtonian limit, we want the action to be also invariant. \\[0.2cm]
Till now, we know three invariant (scalar) quantitites: $m, c, d\tau (\mathrm{or\ } ds)$. Let us construct a quantity with same dimension as of Lagrangian. We then define:
$$\Ss_{\mathrm{rel}} = \int \dd\tau\  L := \SA\int \dd\tau \ mc^2$$
where we replaced the Lagrangian by $\SA mc^2$, $\SA$ is just some constant. This is just our guess but let us see what happens in the Newtonian limit! Note that, using Eq. \ref{eqn:prop}, we can write the action as:
\begin{align*}
    \Ss_{\mathrm{rel}} &= \SA mc^2 \int \dd t \ \sqrt{1- \sfrac{v^2}{c^2}}\\
    &=\SA mc^2\int \dd t \ \qty(1 - \sfrac{v^2}{2c^2}+\SO(\sfrac{v^4}{c^4}))\\
    &=-\SA \int \dd t\ \qty(\frac{1}{2}mv^2 - mc^2)
\end{align*}
Well, this looks like kind of $T-V$ if we identity the rest energy as some `potential' energy, provided $\SA=-1$. Thus, in the Newtonian limit atleast, this defnition of \textit{relativistic action} makes sense kinda. Hence we have:
\begin{equation}\label{eqn:action}
    \Ss_{\mathrm{rel}} := \int \dd \tau \ (-mc^2)
\end{equation}
Note that, we had previously seen that, $\norm{\fvec{u}}^2 = -c^2$ and hence we can replace $-c$ with this quantity in the action expression. Using this we obtain:

$$\Ss_{\mathrm{rel}} = \int \dd \tau \ m \norm{\fvec{u}}^2 = \int \dd \tau \ m\ g_{\mu\nu}u^\mu u^\nu$$
Now, we know that we can derive the ELEOM from the least action principle ($\delta \Ss =0$). We have:
$$\pdv{L}{x^\lambda} =m \partial_\lambda g_{\mu\nu}\ u^\mu u^\nu \qquad \qquad \pdv{L}{u^\lambda} = 2mg_{\lambda\nu}u^\nu$$

Now, the metric could be a function of the coordinates and hence we have:
$$\dv{\tau}\qty(\pdv{L}{u^\lambda}) = 2m\times \qty(g_{\lambda\nu}\dv{u^\nu}{\tau}+\partial_\sigma g_{\lambda\nu}u^\sigma u^\nu)$$
Plugging this into the ELEOM, we obtain:
\begin{align*}
    2m\times \qty(g_{\lambda\nu}\dv{u^\nu}{\tau}+\partial_\sigma g_{\lambda\nu}u^\sigma u^\nu) &=m \partial_\lambda g_{\mu\nu}\ u^\mu u^\nu\implies g_{\lambda\nu}\dv{u^\nu}{\tau} &= \frac{1}{2}\qty(\partial_\lambda g_{\mu\nu}\ u^\mu u^\nu) -\partial_\sigma g_{\lambda\nu}u^\sigma u^\nu
\end{align*}
Then, contracting the metric on the left side with $g^{\rho\lambda}$, we get:
$$\dv{u^\rho}{\tau} = g^{\lambda \rho}\times \qty(\frac{1}{2}\partial_{\lambda}g_{\sigma\nu}-\partial_{\sigma}g_{\lambda\nu})u^{\sigma}u^\nu$$
Let us manipulate the term in the bracket:
\begin{align*}
    &\quad \qty(\frac{1}{2}\partial_{\lambda}g_{\sigma\nu}-\partial_{\sigma}g_{\lambda\nu})u^{\sigma}u^\nu\\
    &=\frac{1}{2}\partial_{\lambda}g_{\sigma\nu}u^{\sigma}u^\nu - \frac{1}{2}\qty(\partial_{\sigma}g_{\lambda\nu}+\partial_{\sigma}g_{\nu\lambda})u^{\sigma}u^\nu\\
    &=\frac{1}{2}\partial_{\lambda}g_{\sigma\nu}u^{\sigma}u^\nu - \frac{1}{2}\partial_{\sigma}g_{\lambda\nu}u^{\sigma}u^\nu-\frac{1}{2}\partial_{\sigma}g_{\nu\lambda}u^{\sigma}u^\nu \\
    &=\frac{1}{2}\partial_{\lambda}g_{\sigma\nu}u^{\sigma}u^\nu - \frac{1}{2}\partial_{\sigma}g_{\lambda\nu}u^{\sigma}u^\nu-\frac{1}{2}\partial_{\nu}g_{\sigma\lambda}u^{\nu}u^\sigma\\
    &=\frac{1}{2}\qty(\partial_{\lambda}g_{\sigma\nu} - \partial_{\sigma}g_{\lambda\nu}-\partial_{\nu}g_{\sigma\lambda})u^{\sigma}u^\nu
\end{align*}
Above, we had used the symmetry of the metric tensor and then we renamed the dummy indices in the sum. We substitute this in the previous equation and obtain:
\begin{align*}
    \dv{u^\rho}{\tau} +\frac{1}{2}g^{\lambda \rho}\qty( \partial_{\sigma}g_{\lambda\nu}+\partial_{\nu}g_{\sigma\lambda} -\partial_{\lambda}g_{\sigma\nu})u^{\sigma}u^\nu=0\implies \dv{u^\rho}{\tau} + \chris{\rho}{\sigma}{\nu}u^\sigma u^\nu = 0
\end{align*}
This is exactly the \textit{geodesic} equation that we had obtained earlier. Hence, using this definition of the action, we can recover the geodesic equation. Thus, \textbf{geodesics are paths of extreme action!}